\documentclass[11pt,a4paper]{article}

% ----- Packages -----
\usepackage[utf8]{inputenc}
\usepackage[T1]{fontenc}
\usepackage{lmodern}
\usepackage[margin=2.2cm]{geometry}
\usepackage{amsmath,amssymb}
\usepackage{siunitx}
\usepackage{booktabs}
\usepackage{multirow}
\usepackage{array}
\usepackage{graphicx}
\usepackage{listings}
\usepackage{xcolor}
\usepackage{hyperref}
\usepackage{enumitem}
\usepackage{tcolorbox}
\usepackage{fancyhdr}
\usepackage{titlesec}
\usepackage{lastpage}

% ----- Hyperref -----
\hypersetup{
  colorlinks=true,
  linkcolor=blue!50!black,
  urlcolor=blue!60!black,
  citecolor=teal!50!black
}

% ----- Colored boxes -----
\tcbuselibrary{skins,breakable}

\newtcolorbox{stepbox}[2][]{
  colback=blue!4!white, colframe=blue!55!black,
  fonttitle=\bfseries, coltitle=white, colbacktitle=blue!55!black,
  title={#2}, arc=2mm, boxrule=0.5pt,
  left=8pt, right=8pt, top=4pt, bottom=4pt, breakable, #1
}

\newtcolorbox{stopcheck}[1][]{
  colback=orange!8!white, colframe=orange!75!black,
  fonttitle=\bfseries, coltitle=white, colbacktitle=orange!75!black,
  title={STOP \& CHECK}, arc=2mm, boxrule=0.5pt,
  left=8pt, right=8pt, top=4pt, bottom=4pt, breakable, #1
}

\newtcolorbox{warningbox}[1][]{
  colback=red!6!white, colframe=red!70!black,
  fonttitle=\bfseries, coltitle=white, colbacktitle=red!70!black,
  title={SAFETY / DATA INTEGRITY}, arc=2mm, boxrule=0.5pt,
  left=8pt, right=8pt, top=4pt, bottom=4pt, breakable, #1
}

\newtcolorbox{tipbox}[1][]{
  colback=green!5!white, colframe=green!50!black,
  fonttitle=\bfseries, coltitle=white, colbacktitle=green!50!black,
  title={TIP}, arc=2mm, boxrule=0.5pt,
  left=8pt, right=8pt, top=4pt, bottom=4pt, breakable, #1
}

\newtcolorbox{keybox}[1][]{
  colback=purple!4!white, colframe=purple!55!black,
  fonttitle=\bfseries, coltitle=white, colbacktitle=purple!55!black,
  title={KEY IDEA}, arc=2mm, boxrule=0.5pt,
  left=8pt, right=8pt, top=4pt, bottom=4pt, breakable, #1
}

% ----- Headers -----
\pagestyle{fancy}
\fancyhf{}
\rhead{\small SOP v4.0 -- DIW Printing Parameters (per-shape suite)}
\lhead{\small BioPoli - PEMM/COPPE/UFRJ}
\rfoot{\small \thepage\ / \pageref{LastPage}}
\renewcommand{\headrulewidth}{0.4pt}

% ----- Code style -----
\lstdefinestyle{shell}{
  basicstyle=\ttfamily\footnotesize,
  backgroundcolor=\color{gray!8},
  frame=single, rulecolor=\color{gray!40},
  columns=fullflexible, breaklines=true,
  showstringspaces=false, tabsize=2
}

% ----- Title spacing -----
\titlespacing*{\section}{0pt}{12pt}{6pt}
\titlespacing*{\subsection}{0pt}{8pt}{4pt}

\title{\textbf{Standard Operating Procedure (SOP) v4.0:\\
DIW 3D Printing Calibration via a Per-Shape Test Suite}\\[4pt]
\large \textsl{One STL per benchmark geometry, swept over
$(v_\mathrm{print}, \mathrm{EM})$ at fixed slicer settings, paired
with a flash$+$no-flash iPhone photograph for downstream
image-based scoring.}}
\author{BioPoli - PEMM/COPPE/UFRJ \\[2pt]
\small Document owner: Thiago~M.~C.~Rodrigues \quad
Version: 4.0 \quad Date: \today}
\date{}

\begin{document}
\maketitle
\thispagestyle{fancy}

\begin{tcolorbox}[colback=blue!3!white,colframe=blue!50!black,
                  arc=2mm,boxrule=0.4pt,left=8pt,right=8pt]
\textbf{What changed from v3.0.}\\
v3.0 collapsed the C1--C4 calibration battery into a single
\textbf{Consolidated Calibration Coupon (CCC\_v1)}. v4.0 returns to
a \emph{per-shape benchmark suite}: each printability quantity is
measured on its own purpose-built specimen, swept over a
$(v_\mathrm{print}, \mathrm{EM})$ grid. The shift is motivated by:
\begin{itemize}[leftmargin=1.4em,itemsep=1pt]
\item \textbf{Reviewer-familiar geometries.} Hollow cubes, towers,
0/90 scaffolds, raster (S-test) and bridges are the standard
benchmarks in the DIW literature. Reporting on these directly is
defensible without explaining a one-off coupon design.
\item \textbf{Quantitative image analysis at the shape scale.}
Each shape isolates one failure mode, which simplifies the
image-processing pipeline (Section~\ref{sec:imaging}) and gives a
single \emph{printability value} per shape per operating point.
\item \textbf{Direct slicer authoring.} G-codes are sliced
once per shape in PrusaSlicer and reused across inks. EM and
$v_\mathrm{print}$ are swept by re-slicing, not by hand-baking $E$
values.
\end{itemize}

\textbf{Photography protocol (new in v4).} Two iPhone photographs
per G-code, taken from the same mounted position --- one with the
built-in flash, one without. Substrate background is either
\textbf{black} or \textbf{blue} matte material (subfoldered).
\end{tcolorbox}

\tableofcontents
\bigskip

% =======================================================
\section{Why a per-shape suite}

\begin{keybox}
A single geometry cannot simultaneously isolate line width,
infill quality, layer-on-layer stacking, vertical stability,
overhang behaviour, and pore fidelity. Each of these failure
modes \emph{is} a shape. The community-standard benchmark
geometries (hollow cube, raster/S, scaffold 0/90, tower, bridge)
were designed to do exactly that.
\end{keybox}

CCC\_v1 (v3.0) prioritised throughput --- one specimen per
operating point --- at the cost of measuring four heterogeneous
quantities on the same print. v4.0 reverses that trade. Each
benchmark is sliced once, swept over $(v_\mathrm{print},
\mathrm{EM})$, photographed under controlled lighting, and scored
by a per-shape metric. The cost is roughly $5\times$ the number of
prints, offset by:
\begin{itemize}[leftmargin=1.4em,itemsep=2pt]
\item Direct comparison against published bioprinting benchmarks.
\item One photograph $\to$ one quantitative score per shape,
through the imaging pipeline of Section~\ref{sec:imaging}.
\item No bespoke CAD --- all geometries are public-domain
benchmarks already used in the field.
\end{itemize}

% =======================================================
\section{Shape inventory}
\label{sec:shapes}

\begin{table}[h]
\centering
\small
\begin{tabular}{p{2.6cm} p{4.4cm} p{4.6cm} p{3.4cm}}
\toprule
Shape & STL location & Quantifies & Failure modes \\
\midrule
Hollow cube & \texttt{Cubes/cubes.stl} &
  Wall fidelity, corner sharpness, perimeter continuity &
  Wall collapse, corner rounding \\
Raster / S-test & \texttt{S-Test/Teste de Printabilidade Raster.stl} &
  Printability index (line straightness and uniformity in a raster
  pattern) &
  Under/over-extrusion, stringing \\
Area cubes (full-fill, varying height) &
  \texttt{Area-Cube/Modelo Area Cubos.stl} &
  Infill density, layer adhesion, top-surface flatness &
  Inter-line gaps, top-surface roughness \\
Tower 20\,mm & \texttt{Tower-20mm/} &
  Short-tower stability, base bulge &
  Buckling onset, base spreading \\
Tower 40\,mm & \texttt{Tower-40mm/} &
  Tall-tower buckling limit &
  Lateral collapse, base creep \\
Scaffold 0/90 & \texttt{Scaffold 0-90/} &
  Pore fidelity, layer-on-layer registration &
  Pore closure, strand sag into pore \\
Bridge / suporte & \texttt{Bridge/Modelo Suporte.stl} &
  Bridging / overhang &
  Mid-span sag, dropout \\
Misc (Pot, Star, Pyramid, Maze) & \texttt{Misc/} &
  Complex-geometry validation &
  Geometry-specific (e.g.\ thin-feature failure) \\
\bottomrule
\end{tabular}
\caption{Benchmark shape inventory. STLs sliced once in PrusaSlicer,
swept across $(v_\mathrm{print}, \mathrm{EM})$. ``Failure modes''
indicate the dominant artefact each shape is designed to expose.}
\label{tab:shapes}
\end{table}

\subsection{Folder organisation under \texttt{Impressao 3D/}}

\begin{lstlisting}[style=shell]
Impressao 3D/
  Area-Cube/      <- STL + 15 gcodes (Vp x Fr sweep)
  Bridge/         <- STL only
  Cubes/          <- STL only
  Misc/Pot/       <- STL + gcodes (in progress)
  Misc/Star/      <- STL + gcodes (in progress)
  S-Test/         <- STL + 20 gcodes (Vp x Flow sweep)
    SF5.5/
      Black/      <- iPhone photos on BLACK substrate
      Blue/       <- iPhone photos on BLUE substrate
  Scaffold 0-90/  <- STL + gcodes (in progress)
  Tower-20mm/     <- STL + gcodes (in progress)
  Tower-40mm/     <- STL + gcodes (in progress)
  OLD/            <- CCC_v1 archival (do not use)
\end{lstlisting}

The \texttt{SF5.5/} subfolder is an ink identifier (here:
C15 + sunflower-oil nanoemulsion @ 5.5\%). \texttt{Black/} and
\texttt{Blue/} are \emph{substrate background colours}, not ink
colours. The same operating point is photographed against both
backgrounds when contrast against one background is degenerate
(e.g.\ white ink fades against a high-key flash on blue but reads
on black). Subsequent inks introduce a sibling subfolder
(e.g.\ \texttt{C15/}, \texttt{Bozano/}) at the same level as
\texttt{SF5.5/}.

% =======================================================
\section{Naming convention for G-codes}
\label{sec:naming}

\begin{table}[h]
\centering
\small
\begin{tabular}{p{3.4cm} p{4.0cm} p{6.6cm}}
\toprule
Token & Meaning & Sweep values currently in folder \\
\midrule
\texttt{Vp$N$} &
$v_\mathrm{print}$ (print-head speed) in mm/s.
\textbf{Not} piston velocity. &
$N \in \{5, 7, 10, 15, 30\}$ \\
\texttt{Fr$X$} or \texttt{Flow$X$} &
Extrusion Multiplier (the PrusaSlicer field). $1.0$ is
no-correction. &
Area-Cube: $\{1, 1.5, 2\}$; \quad
S-Test: $\{1, 1.57, 2, 2.5\}$. The value $1.57$ matches
the Cross-model prediction $1/k_\mathrm{flow}$ for SF5.5
(NE-loaded CMC). \\
\bottomrule
\end{tabular}
\caption{Filename tokens. Example: \texttt{Vp10-Fr1,5.gcode}
$\Rightarrow$ $v_\mathrm{print}=10$ mm/s, $\mathrm{EM}=1.5$.}
\label{tab:naming}
\end{table}

\begin{warningbox}
\textbf{``Vp'' here is print-head speed, not piston velocity.}
In the rheology / simulation pipeline (CLAUDE.md, MATLAB
\texttt{run\_solver\_v3.m}), $V_p$ denotes \emph{piston advance
velocity} (typically $\sim 0.01$ mm/s for DIW). The G-code
filenames in \texttt{Impressao 3D/} use \texttt{Vp} as shorthand
for $v_\mathrm{print}$, the slicer's perimeter speed (5--30~mm/s).
Both quantities are physically distinct and orders of magnitude
apart. Always read the gcode's \texttt{F$\,\langle\mathrm{value}\rangle$}
feedrate to confirm: \texttt{F600} = 10~mm/s = \texttt{Vp10}.
\end{warningbox}

% =======================================================
\section{PrusaSlicer configuration}
\label{sec:slicer}

The fixed-field profile from SOP v3.0 carries over unchanged
(filament diameter 14.3~mm, retract 0, fan off, extrusion width
0.5~mm, etc.). Only the variable fields differ between G-codes
within a sweep.

\subsection{Fixed fields (one-time printer profile)}

\begin{table}[h]
\centering
\small
\begin{tabular}{p{5.2cm} p{3.4cm} p{6.2cm}}
\toprule
Slicer field & Value & Why \\
\midrule
Filament diameter & $\mathbf{14.3\,\mathrm{mm}}$ &
  $= 2 R_\mathrm{syringe}$, BD 10 mL. Makes 1 mm slicer-$E$ = 1 mm
  piston. Critical. \\
Filament density & $1.00\,\mathrm{g/cm^3}$ & Bookkeeping. \\
Cooling fan & OFF & DIW. \\
Retract length & 0 mm & DIW. \\
Retract Z-hop & 0 mm & DIW. \\
Travel speed & 60 mm/s & Avoid head-drag through wet ink. \\
Skirt loops & 1 & Prime line. \\
Brim & 0 & DIW. \\
Top/bottom solid layers & 0 (except Area-Cube top: 2) &
  Allow top-surface flatness reading on Area-Cube only. \\
Perimeters & 1 (Hollow Cube, Scaffold, Tower) or 2 (Bridge) &
  Single perimeter for transparent failure-mode reading. \\
Extrusion width & $\mathbf{0.5\,\mathrm{mm}}$ &
  Matches 21G needle ID (0.515 mm). \\
First-layer extrusion width & $0.5\,\mathrm{mm}$ & No squish. \\
First-layer speed & $=$ Perimeter speed &
  Do not let PrusaSlicer slow the first layer. \\
\bottomrule
\end{tabular}
\caption{Fixed PrusaSlicer fields. Save in a printer + filament
profile pair (e.g.\ \texttt{DIW\_BD10mL\_21G}).}
\label{tab:slicer_fixed}
\end{table}

\subsection{Variable fields (per G-code)}

\begin{table}[h]
\centering
\small
\begin{tabular}{p{4.2cm} p{3.8cm} p{6.4cm}}
\toprule
Slicer field & Set to & Source / formula \\
\midrule
Filament $\to$ Extrusion Multiplier &
  $\mathrm{EM} = 1/k_\mathrm{flow,pred}$ &
  From \texttt{slicer\_lookup\_*.csv} (output of
  \texttt{run\_solver\_v3.m}). For SF5.5: 1.57. For an unknown ink,
  start at 1.0 and refine. \\
Filament $\to$ Max volumetric speed &
  $Q_\mathrm{cap} = \pi R_\mathrm{syr}^2 V_p$ &
  For BD 10 mL: $Q_\mathrm{cap}\,[\mathrm{mm^3/s}] = 160.6 \cdot
  V_\mathrm{p,piston}\,[\mathrm{mm/s}]$. At $V_\mathrm{p,piston}
  = 0.01$ mm/s, $Q_\mathrm{cap} = 1.606$ mm$^3$/s. \\
Print $\to$ Perimeter speed (and First-layer speed) &
  $v_\mathrm{print}$ from the test matrix &
  Sweep values: 5, 7, 10, 15, 30 mm/s (the ``Vp'' filename token). \\
\bottomrule
\end{tabular}
\caption{Per-G-code variable fields. All other parameters in
Table~\ref{tab:slicer_fixed} remain fixed.}
\label{tab:slicer_variable}
\end{table}

\subsection{Predicted EM (current ink panel)}

\begin{table}[h]
\centering
\small
\begin{tabular}{lccc}
\toprule
Ink & $k_\mathrm{flow,pred}$ & $\mathrm{EM} = 1/k_\mathrm{flow}$ & Notes \\
\midrule
C15 & 0.83 & \textbf{1.20} & Baseline CMC \\
C15 Gira 5.5 (SF5.5) & 0.64 & \textbf{1.57} &
  Nanoemulsion-loaded, lowest recovery \\
Bozano & 1.06 & \textbf{0.94} & Commercial reference \\
\bottomrule
\end{tabular}
\caption{Cross-model predictions, $V_\mathrm{p,piston}=0.01$ mm/s,
21G. The 1.57 value appears in the S-Test sweep as the
SF5.5-tuned operating point.}
\end{table}

\subsection{Start / End G-code}

PrusaSlicer-default homing is preserved in v4 (the manual
\texttt{M0} pause used in v3 is no longer required because each
shape is run as a self-contained print). Confirm the
slicer-generated header begins with \texttt{M92 E2651.19} and
\texttt{G28}; the head will return to the bed corner for the
first move.

\begin{lstlisting}[style=shell]
; Start G-code (PrusaSlicer default for DIW printer)
M92 E2651.19
G28          ; home all axes
G1 Z5 F5000  ; lift nozzle
G21          ; mm
G90          ; absolute coords
M82          ; absolute E
G92 E0
G1 E1 F3000  ; prime
G92 E0

; End G-code
G1 Z30 F600
M84
\end{lstlisting}

\begin{warningbox}
If the printer's mechanical Z-zero is set above the substrate
(safety lift) rather than at substrate contact, jog Z down and
issue \texttt{G92 Z0} \emph{before} starting the print --- the
slicer assumes substrate contact at $Z=0$.
\end{warningbox}

% =======================================================
\section{Test matrix and execution order}
\label{sec:matrix}

\subsection{Sweep grids}

\begin{table}[h]
\centering
\small
\begin{tabular}{p{2.8cm} p{2.8cm} p{2.8cm} p{2.8cm}}
\toprule
Shape & $v_\mathrm{print}$ (mm/s) & EM & Total G-codes \\
\midrule
Area-Cube & 5, 7, 10, 15, 30 & 1.0, 1.5, 2.0 & 15 \\
S-Test    & 5, 7, 10, 15, 30 & 1.0, 1.57, 2.0, 2.5 & 20 \\
Cubes (hollow) & 5, 7, 10, 15, 30 & 1.0, 1.57, 2.0 & 15 (planned) \\
Tower-20mm & 5, 10, 15 & 1.0, 1.57 & 6 (planned) \\
Tower-40mm & 5, 10 & 1.0, 1.57 & 4 (planned) \\
Scaffold 0/90 & 5, 7, 10, 15, 30 & 1.0, 1.57, 2.0 & 15 (planned) \\
Bridge & 5, 10, 15 & 1.0, 1.57 & 6 (planned) \\
Misc & per-shape (single best-EM scout) & 1.57 (or refined) & TBD \\
\bottomrule
\end{tabular}
\caption{Sweep grids per shape. ``Planned'' entries are
under generation by the operator; the SOP is forward-compatible.
The S-Test grid is the most exhaustive because the raster
printability index is the headline metric in the literature.}
\label{tab:sweep_grids}
\end{table}

\subsection{Recommended execution order}

\begin{enumerate}[leftmargin=1.4em,itemsep=2pt]
\item \textbf{S-Test full grid (20 prints).} Establishes the
$(v_\mathrm{print}, \mathrm{EM})$ window for the ink. The raster
geometry is the most sensitive to misalignment of EM with the
ink's recovery, so it converges first.
\item \textbf{Area-Cube grid (15 prints).} Validates the
S-Test-derived window for filled regions.
\item \textbf{Hollow Cube + Scaffold + Tower + Bridge} at the
two or three operating points that survived S-Test scoring.
\item \textbf{Misc shapes} as final demonstration prints at the
single best operating point.
\end{enumerate}

% =======================================================
\section{Photography protocol}
\label{sec:imaging}

Every printed specimen is photographed \emph{immediately after
the print finishes} (before any relaxation) and \emph{again at
$t=30$ min relaxation} (for shapes where slump matters: Tower,
Scaffold, Bridge, Area-Cube). The intermediate window is not
photographed.

\subsection{Per print: two iPhone images}

\begin{stepbox}{Image pair acquisition}
\begin{enumerate}[leftmargin=1.4em,itemsep=2pt]
\item Mount iPhone in a fixed clamp directly above the print bed,
camera plane parallel to the substrate. Mark the mount position
on the bench --- it must be reproducible across all prints in a
sweep.
\item Set the iPhone to \textbf{standard Photo mode} (not Portrait,
not Live, not RAW unless ProRAW is consistently used). HDR may
remain on \emph{Auto} but must be consistent within a sweep.
\item Capture \textbf{image A: with flash}. Tap to focus on the
specimen centre, do not adjust exposure compensation.
\item Capture \textbf{image B: without flash}, identical framing,
same focus point.
\item Filename convention:
\texttt{<shape>\_Vp<N>\_Fr<X>\_<ink>\_<bg>\_<flash|noflash>.jpg}
\begin{itemize}[leftmargin=1.4em]
\item Example:
\texttt{S-Test\_Vp10\_Flow1,57\_SF5.5\_Black\_flash.jpg}
\end{itemize}
\item Save the pair into the corresponding
\texttt{Impressao 3D/<shape>/<ink>/<bg>/} subfolder.
\end{enumerate}
\end{stepbox}

\begin{keybox}
\textbf{Why a flash $+$ no-flash pair.} The flash image
\emph{kills shadows and saturates highlights}, exposing geometry
boundaries with high contrast --- ideal for edge detection. The
no-flash image \emph{preserves the natural reflection / opacity
gradient}, exposing layer-to-layer surface defects --- ideal for
texture analysis. Together the pair carries more information
than a single ``optimised'' shot, and an automated pipeline can
choose channel-by-channel which image to use.
\end{keybox}

\subsection{Framing and substrate}

\begin{table}[h]
\centering
\small
\begin{tabular}{p{4.2cm} p{10.2cm}}
\toprule
Aspect & Specification \\
\midrule
Camera axis & Perpendicular to substrate (top-down) for Cubes,
  Area-Cube, S-Test, Scaffold; perpendicular to tower axis
  (side-on) for Tower-20/40 and Bridge. \\
Working distance &
  Fixed; record it on the run sheet. The image pipeline cannot
  recover absolute scale without either (a) a reference object in
  frame or (b) a recorded working distance + iPhone lens focal
  length. \\
In-frame scale reference &
  \emph{Strongly recommended}: place a printed RefBar (10$\times$1
  mm, 0.4 mm tall) or a ruler segment alongside every specimen.
  This is the same RefBar used in v3.0; STL still in
  \texttt{OLD/CCC\_v1.stl}. \\
Background &
  Matte black or matte blue card, large enough to fill the field
  of view. Avoid glossy substrates --- the flash returns a
  specular peak that corrupts edge detection. \\
Ambient lighting &
  Consistent within a sweep. Overhead fluorescent OK, mixed
  daylight + tungsten not OK. Record on the run sheet. \\
\bottomrule
\end{tabular}
\caption{Framing rules. Working distance and lighting are
recorded because the image pipeline (Section~\ref{sec:imaging_algo})
operates in pixel space until a reference object provides
mm-per-pixel calibration.}
\label{tab:framing}
\end{table}

% =======================================================
\section{Image-analysis pipeline (in development)}
\label{sec:imaging_algo}

The pipeline ingests the flash + no-flash pair and emits a
single \emph{printability value} per specimen, in pixel units
until a scale reference is provided. The full algorithm is
under development; this section records the design intent so
that the photography protocol remains forward-compatible.

\subsection{Inputs and outputs}

\begin{itemize}[leftmargin=1.4em,itemsep=2pt]
\item \textbf{Inputs.} Two JPEGs per print:
\texttt{<...>\_flash.jpg} and \texttt{<...>\_noflash.jpg}, plus
the gcode-derived metadata $(shape, v_\mathrm{print},
\mathrm{EM}, h_\mathrm{layer}, ink, bg)$.
\item \textbf{Outputs (pixel domain).} Per shape, a vector of
geometric defect metrics --- e.g.\ for the S-Test: line-width
mean and SD, raster spacing mean and SD, fraction of expected
raster nodes resolved. For the hollow cube: perimeter
deviation (px), corner radius (px). For towers: base-bulge
width vs nominal (px). Pixel $\to$ mm conversion is deferred
to a calibration revision the operator will supply.
\item \textbf{Headline scalar.} Per shape, a single
\emph{printability value} $\in [0, 1]$ combining the
shape-specific metrics. Aggregation rule TBD during the
brainstorming phase.
\end{itemize}

\subsection{Stages (intent, not implementation)}

\begin{enumerate}[leftmargin=1.4em,itemsep=2pt]
\item \textbf{Pair registration.} Align the flash and no-flash
images pixel-wise (small mount drift between shots is expected).
\item \textbf{Background segmentation.} Background colour (black
or blue) is metadata; threshold-based segmentation produces a
foreground mask.
\item \textbf{Reference-object detection.} If a RefBar or ruler
is in frame, extract a px/mm scale; otherwise return pixel-only
metrics.
\item \textbf{Per-shape metric extraction.} A shape-specific
module (one per row of Table~\ref{tab:shapes}) computes the
defect metrics.
\item \textbf{Score aggregation.} Per-shape printability value
$\in [0,1]$.
\end{enumerate}

% =======================================================
\section{Measurement protocols (manual fallback)}
\label{sec:measure}

Until the image-analysis pipeline is operational, each shape is
scored manually. The protocols below are the human counterparts
of the automated metrics in Section~\ref{sec:imaging_algo}.

\subsection{Hollow cube}
Visually rate perimeter continuity (closed vs broken) and
corner sharpness (sharp / rounded / collapsed) on a 0--3 scale.
Photograph top-down; measure inner and outer edge length with
caliper; deviation from nominal is the headline number.

\subsection{S-Test / raster}
Photograph top-down. Count the number of expected raster
turn-around nodes that are resolved (i.e.\ where the strand
clearly reverses without fusing into the previous pass).
Printability index $= N_\mathrm{resolved} / N_\mathrm{expected}$.

\subsection{Area cubes (different heights)}
For each height, measure the top-surface gap fraction
(fraction of the patch where adjacent raster lines have not
fused) on a top-down photograph. Caliper the total stack
height vs $N_\mathrm{sliced} \cdot h_\mathrm{layer}$ for the
stacking ratio (carries over from v3.0 C3).

\subsection{Tower 20\,mm and 40\,mm}
Side-on photograph during and after the print. Record the
layer count at which base width first exceeds nominal by 10\%
($N_\mathrm{exp}$). Compare to the SAOS prediction in
\texttt{printing\_parameters\_per\_ink.csv}.

\subsection{Scaffold 0/90}
Top-down photograph after relaxation. Count the fraction of
pores that remain open (no strand sag bridging the pore). Pore
closure rate is the headline metric.

\subsection{Bridge}
Side-on photograph. Measure mid-span sag distance with caliper
against a known span length. Sag fraction $=$ sag / span.

\subsection{Misc shapes}
Single qualitative verdict per shape: \emph{pass} (geometry
recognisable, no major defects) / \emph{partial} (recognisable
but with one major defect) / \emph{fail}.

% =======================================================
\section{Run sheet}
\label{sec:runsheet}

\begin{center}
\small
\begin{tabular}{|p{0.95\linewidth}|}
\hline
\textbf{Per-Shape Print Run Sheet} \quad
Date: \rule{2.5cm}{0.4pt} \quad
Operator: \rule{3cm}{0.4pt} \\
\hline
\textbf{Print metadata} \\
Shape: \rule{4cm}{0.4pt} \quad
G-code file: \rule{6cm}{0.4pt} \\
$v_\mathrm{print}$ = \rule{1.4cm}{0.4pt}\,mm/s \quad
EM = \rule{1.2cm}{0.4pt} \quad
$h_\mathrm{layer}$ = \rule{1.4cm}{0.4pt}\,mm \\
Ink: \rule{4cm}{0.4pt} \quad
Batch: \rule{2.5cm}{0.4pt} \quad
Needle: $\square$ 21G $\;\square$ 22G \\
Substrate background: $\square$ Black $\;\square$ Blue $\;\square$ \rule{2.5cm}{0.4pt} \\[3pt]
\hline
\textbf{Pre-print} \\
Filament diameter set to 14.3 mm? $\square$ Yes $\;\square$ No (STOP) \\
Z-zero set at substrate contact? $\square$ Yes \\
\hline
\textbf{Photography (at end of print)} \\
Mount position reproduced? $\square$ Yes \\
Working distance: \rule{2cm}{0.4pt}\,cm \quad
Ambient lighting: \rule{4cm}{0.4pt} \\
Image A (flash) saved? $\square$ \quad
Image B (no flash) saved? $\square$ \\
$t=30$ min repeat pair (if applicable)? $\square$ \\
\hline
\textbf{Manual verdict (until imaging pipeline live)} \\
\rule{0.93\linewidth}{0.4pt} \\
\rule{0.93\linewidth}{0.4pt} \\
\hline
\end{tabular}
\end{center}

% =======================================================
\section{Deprecation notice}

The Consolidated Calibration Coupon workflow described in
SOP v3.0 (\texttt{Latex/Printing\_Parameters\_SOP\_v3\_EN.tex}
and \texttt{CCC\_v1.stl}) is \textbf{deprecated as of
\today}. v3.0 and the legacy G-codes under
\texttt{Impressao 3D/OLD/} remain in the repository for
archival reference. The active workflow for all new prints is
the per-shape suite defined in this document.

% =======================================================
\section{Changelog}

\begin{description}
\item[v4.0 (\today).] Replaced CCC\_v1 single-coupon workflow with
per-shape benchmark suite (hollow cube, area cubes, S-test,
tower 20/40, scaffold 0/90, bridge, misc). New iPhone flash $+$
no-flash photography protocol with shape/ink/background
subfolder hierarchy. Sweep grids over $(v_\mathrm{print},
\mathrm{EM})$ documented per shape. Image-analysis pipeline
placeholder added pending algorithm finalisation.
\item[v3.0.] CCC\_v1 consolidated coupon with C1--C4 metrics on one specimen.
\item[v2.0.] 18 hand-coded G-codes per ink (C1--C4 battery).
\end{description}

\end{document}
